% =========================================================================
% Plan de gestión de riesgos — versión 1 (emitida el 26-07-2026, cierre de S1).
% Documento independiente de la memoria: la v0 es el cap. 7 y no se toca
% (línea base); cada versión nueva es un documento completo (ver README.md).
% REGLA DE CONTENIDO: las tablas reproducen LITERALMENTE las de la v0
% (cap. 7), salvo los riesgos cuyo cambio documenta el seguimiento de S1
% (R4: cláusula de mitigación eliminada; R10: estado Resuelto). Diferencias
% de forma respecto al capítulo: sin \ref (se escribe el número de capítulo)
% y sin \ac (el acrónimo directamente), porque este documento compila solo.
% Diseño heredado del plan de riesgos de Gestión de Proyectos del autor,
% adaptado a la configuración global de la memoria (preámbulo compartido).
% Compilar DESDE ESTA CARPETA (2 pases, sin biber):
%   pdflatex --enable-installer -interaction=nonstopmode -file-line-error
%     -output-directory=build plan-riesgos-v1-26-07-2026.tex
% =========================================================================
\documentclass[11pt,a4paper]{article}

% =========================================================================
% Preámbulo COMPARTIDO de los documentos independientes de riesgos
% (docs/thesis/plan-riesgos/ y docs/thesis/seguimiento-riesgos/).
%
% Estos documentos NO forman parte de main.tex (la memoria): son entregables
% sueltos en clase article, con el diseño heredado de los .tex de Gestión de
% Proyectos del autor, pero adaptados a la configuración global de la
% memoria (Carlito, márgenes 2,5 cm, parskip 6 pt, títulos de la guía).
% Un solo preámbulo para toda la familia: un ajuste de estilo aquí se
% propaga a todos los seguimientos futuros (uno por sprint hasta S37).
%
% RUTAS: los \input son relativos a la carpeta del documento, no a esta;
% compilar SIEMPRE desde la carpeta del .tex (plan-riesgos/ o
% seguimiento-riesgos/):
%   pdflatex --enable-installer -interaction=nonstopmode -file-line-error
%     -output-directory=build <documento>.tex        (2 pases; sin biber)
% =========================================================================

% --- Codificación (idéntico a packages.tex de la memoria) ---
\usepackage[utf8]{inputenc}
\usepackage[T1]{fontenc}

% --- Idioma: solo español (aquí no hay resumen trilingüe) ---
% es-noshorthands por coherencia con la memoria: acentos en UTF-8, sin
% caracteres activos (el conflicto con biblatex-apa no aplica aquí, pero
% mantener una única configuración de babel evita sorpresas al copiar texto
% entre la memoria y estos documentos).
\usepackage[spanish, es-tabla, es-noshorthands]{babel}

% --- Márgenes: 2,5 cm laterales e inferior, como la memoria ---
% El margen superior se reparte para alojar la cabecera enmarcada del
% diseño de referencia (2 cm de caja + separación) sin comerse el cuerpo:
% el texto arranca a 4 cm del borde y el ancho de texto sigue siendo 16 cm.
\usepackage[a4paper, left=2.5cm, right=2.5cm, bottom=2.5cm,
            top=1.2cm, includehead, headheight=2.2cm, headsep=0.6cm]{geometry}

% --- Tipografía: Carlito, la misma decisión que en style.tex ---
% (clon métrico libre de Calibri; pdflatex no puede usar la real).
% No se replica el arreglo de versalitas de style.tex: aquí no hay
% bibliografía APA y ningún texto usa \textsc.
\usepackage[sfdefault]{carlito}

% --- Espaciado: interlineado sencillo + 6 pt entre párrafos (guía) ---
\usepackage[skip=6pt]{parskip}

% --- Tablas ---
% longtable: el registro y el plan de respuesta superan una página.
% AVISO ESPERADO: "ignored error: Infinite glue shrinkage" en cada salto de
% página de una longtable (bug conocido e inofensivo, latex3/latex2e#1907;
% ver el comentario largo en packages.tex de la memoria).
\usepackage{array}
\usepackage{longtable}
\usepackage{tabularx}
% \RaggedRight (ragged2e) en vez de \raggedright en columnas muy estrechas:
% el \raggedright de LaTeX desactiva de facto la partición silábica y una
% palabra larga («Identificado») desborda una columna de 1,2 cm; la
% variante de ragged2e mantiene el guionado.
\usepackage{ragged2e}
% Las columnas X se componen como m{...}: el contenido de cada celda queda
% centrado verticalmente (aspecto de formulario del diseño de referencia).
\renewcommand{\tabularxcolumn}[1]{m{#1}}
% Tipos de columna del diseño de referencia:
%   L/C = ancho fijo alineado izquierda/centro; Y/Z = X alineada izq./centro.
% En columnas estrechas se prefiere Y (raggedright) a X justificada: la
% justificación en 3-4 cm abre ríos de blancos.
\newcolumntype{L}[1]{>{\raggedright\arraybackslash}m{#1}}
\newcolumntype{C}[1]{>{\centering\arraybackslash}m{#1}}
\newcolumntype{Y}{>{\raggedright\arraybackslash}X}
\newcolumntype{Z}{>{\centering\arraybackslash}X}
% Filas con aire en todas las tablas (el diseño de referencia lo fijaba
% tabla a tabla; aquí una sola vez).
\renewcommand{\arraystretch}{1.25}

% --- Utilidades ---
\usepackage{xcolor}    % \markfield (resaltado) y \todonote (macros.tex)
\usepackage{amssymb}   % \square, \blacksquare, \circlearrowright
\usepackage{xspace}    % lo necesita \bukit (macros.tex)
\usepackage{enumitem}
\setlist[itemize]{leftmargin=*, topsep=2pt, itemsep=2pt}
\setlist[enumerate]{leftmargin=*, topsep=2pt, itemsep=2pt}

% --- Títulos según la guía de estilo, como en style.tex de la memoria ---
% (sin \chapter: la clase article no lo tiene; el resto, idéntico).
\usepackage{titlesec}
\titleformat{\section}
  {\normalfont\bfseries\fontsize{12pt}{14.4pt}\selectfont\raggedright}
  {\thesection.}{0.5em}{\MakeUppercase}
\titleformat{\subsection}
  {\normalfont\bfseries\fontsize{12pt}{14.4pt}\selectfont}
  {\thesubsection.}{0.5em}{}
\titleformat{\subsubsection}
  {\normalfont\itshape\fontsize{12pt}{14.4pt}\selectfont}
  {\thesubsubsection.}{0.5em}{}

% --- Paginación dinámica (regla estricta del autor, 26-07-2026) ---
% PROHIBIDO el salto de página manual también en estos documentos: el
% espacio lo garantizan estos ganchos (o un \Needspace* local ante un bloque
% indivisible), nunca un \newpage. Los .tex de referencia sí usaban
% \newpage; es la única desviación deliberada respecto a ellos.
\usepackage{needspace}
\AddToHook{cmd/section/before}{\Needspace*{5\baselineskip}}
\AddToHook{cmd/subsection/before}{\Needspace*{5\baselineskip}}
\AddToHook{cmd/subsubsection/before}{\Needspace*{4\baselineskip}}
\widowpenalty=10000
\clubpenalty=10000
\displaywidowpenalty=10000
\raggedbottom
% Estos formularios componen mucho texto en columnas de 2-4 cm: se concede
% a TeX un margen de estirado de emergencia antes de desbordar la línea
% (evita Overfull \hbox menores sin tocar el contenido).
\setlength{\emergencystretch}{1em}

% --- Datos globales y macros de la memoria (única fuente: no duplicar) ---
% =========================================================================
% Variables globales del documento.
% Cambiar un dato AQUÍ lo actualiza en todo el documento (portada,
% metadatos del PDF...). No escribir nunca estos datos "a mano" en otros
% archivos: usar siempre estos comandos.
% =========================================================================

% --- Trabajo ---
\newcommand{\thesisTitle}{BUKIT: red social orientada a objetivos}% TODO: título definitivo
\newcommand{\thesisSubtitle}{Memoria del Trabajo de Fin de Grado}%
\newcommand{\thesisDegree}{Grado en Ingeniería Informática de Gestión y Sistemas de Información}%
\newcommand{\thesisCourse}{2026--2027}% curso académico
\newcommand{\thesisDate}{Bilbao, \today}% TODO: <Lugar, día, mes, año> de la entrega/defensa
\newcommand{\thesisYear}{2027}% año de elaboración y defensa

% --- Personas (la plantilla del centro usa "Apellido1 Apellido2, Nombre") ---
\newcommand{\thesisAuthor}{Liviu Deleanu}% nombre natural (para metadatos del PDF)
\newcommand{\thesisAuthorFormal}{Deleanu, Liviu}% TODO: confirmar formato con dos apellidos
\newcommand{\thesisDirector}{Villamañe Gironés, Mikel}%
\newcommand{\thesisCoDirector}{Apellido1 Apellido2, Nombre}% TODO: solo si hay codirector/a

% --- Universidad ---
\newcommand{\thesisUniversity}{Universidad del País Vasco / Euskal Herriko Unibertsitatea}%
\newcommand{\thesisFaculty}{Escuela de Ingeniería de Bilbao}%

% --- Palabras clave del resumen trilingüe ---
\newcommand{\thesisKeywordsEs}{red social, objetivos, monorepo, TypeScript, PostgreSQL}% TODO
\newcommand{\thesisKeywordsEu}{sare soziala, helburuak, monorepo, TypeScript, PostgreSQL}% TODO: revisar euskera
\newcommand{\thesisKeywordsEn}{social network, goals, monorepo, TypeScript, PostgreSQL}% TODO

% =========================================================================
% Comandos propios. Definir aquí cualquier cosa que se repita en el texto:
% si mañana cambia cómo se escribe, se cambia en un único sitio.
% =========================================================================

% Nombre del producto, siempre con el mismo formato.
\newcommand{\bukit}{\textsc{Bukit}\xspace}

% Términos en inglés, en cursiva de forma consistente: \eng{backend}
\newcommand{\eng}[1]{\textit{#1}}

% Notas de trabajo, visibles en rojo en el PDF. Buscar "todonote" y
% eliminarlas todas antes de entregar.
\newcommand{\todonote}[1]{\textcolor{red}{\textbf{[TODO: #1]}}}


% --- Cabecera y pie (diseño de referencia, sin logos) ---
% La caja de cabecera replica la del .tex de referencia (celda ancha de
% contexto + celda estrecha de identidad). La celda derecha llevaba el logo
% del centro: aquí va texto porque los logos UPV/EHU siguen pendientes de
% la fase de diseño (ver style.tex); cuando existan, sustituir por
% \includegraphics. 11 + 4,9 = 15,9 cm < 16 cm de texto: sin desbordes.
\usepackage{fancyhdr}
\newcommand{\riesgosHeader}{%
  \setlength{\tabcolsep}{0pt}%
  \renewcommand{\arraystretch}{1}%
  \begin{tabular}{|p{11cm}|p{4.9cm}|}
    \hline
    \parbox[c][2cm][c]{11cm}{\centering\bfseries
      {\fontsize{16}{19}\selectfont TRABAJO DE FIN DE GRADO\\[4pt]Curso \thesisCourse}%
    } &
    \parbox[c][2cm][c]{4.9cm}{\centering
      {\fontsize{14}{17}\selectfont\bfseries\bukit}\\[3pt]
      {\small\thesisFaculty}%
    }\\
    \hline
  \end{tabular}%
}
\pagestyle{fancy}
\fancyhf{}
\renewcommand{\headrulewidth}{0pt}
\fancyhead[C]{\riesgosHeader}
\fancyfoot[C]{\thepage}
% La primera página usa estilo "plain": se redefine para que también lleve
% la cabecera (mismo truco que el .tex de referencia).
\fancypagestyle{plain}{%
  \fancyhf{}%
  \renewcommand{\headrulewidth}{0pt}%
  \fancyhead[C]{\riesgosHeader}%
  \fancyfoot[C]{\thepage}%
}

% --- Símbolos y campos de formulario del diseño de referencia ---
\newcommand{\fieldline}[1]{\underline{\hspace{#1}}}
\newcommand{\bigline}{\noindent\rule{\textwidth}{0.4pt}}
\newcommand{\checkbox}{\(\square\)}
\newcommand{\checkedbox}{\(\blacksquare\)}
% Símbolo ⟳ = "si se rellena este apartado, nueva versión del plan".
% El original lo resolvía con un SVG externo (paquete svg + Inkscape);
% aquí basta el símbolo de amssymb: cero dependencias nuevas.
\newcommand{\updatesymbol}{\raisebox{-0.2ex}{\(\circlearrowright\)}}
% Resaltado amarillo del identificador de riesgo en los títulos.
\newcommand{\markfield}[1]{\colorbox{yellow!35}{\strut #1}}

% --- Enlaces y metadatos del PDF (al final, como exige hyperref) ---
% El pdftitle lo fija cada documento con \hypersetup tras cargar esto.
\usepackage[hidelinks]{hyperref}
\hypersetup{pdfauthor={\thesisAuthor}}

\hypersetup{pdftitle={Plan de gestión de riesgos — versión 1 (26-07-2026)}}

\begin{document}

\begin{center}
  {\Huge\bfseries\bukit}\\[0.5cm]
  {\Large\bfseries Plan de gestión de riesgos (versión 1)}\\[0.4cm]
  {\large 26 de julio de 2026}
\end{center}
\vspace{0.6cm}

\textbf{Proyecto:} \thesisTitle\ — Trabajo de Fin de Grado de \thesisAuthor\ (curso \thesisCourse).

\textbf{NOTA — cambios respecto a la v0} (cap.~7 de la memoria, versión 0 de 17-07-2026):
\begin{enumerate}
  \item \textbf{R4 (incumplimiento de la normativa de la memoria)}: se elimina de la mitigación la cláusula final «el número de páginas se observa en cada compilación».
  Motivo: el autor excluyó el 26-07-2026 el apartado de extensión (40/60 pp.) de la guía de estilo aplicable a este TFG; sin límite normativo de páginas, vigilarlas ya no mitiga nada y el riesgo queda centrado en la citación (plagio) y el resto de la normativa.
  P, I y severidad no cambian.
  \item \textbf{R10 (indisponibilidad prolongada del autor)}: se materializó en S1 (8 días sin actividad, 18-07 \(\to\) 25-07, por agotamiento laboral y asuntos personales) y queda \textbf{Resuelto} tras aplicar la contingencia en su versión mínima (pausa de facto + re-encaje dentro del propio S1; la historia del sprint se completó el 26-07).
  La evaluación (P1 \(\times\) I3, significativo) se confirma y no cambia.
  \item El resto de riesgos: sin cambios (revisión de S1 en \path{../seguimiento-riesgos/2026-07-26-S1.tex}).
\end{enumerate}

Emitida en la primera retrospectiva del proyecto (cierre de S1).
Registrada en la tabla de seguimiento del cap.~7 de la memoria y en el dev-journal del 26-07-2026.

\section{Identificación, análisis y clasificación de riesgos}

Escalas de valoración (matriz probabilidad--impacto del cap.~7):
\begin{itemize}
  \item \textbf{Probabilidad (P) e impacto (I):} de 1 a 3.
  \item \textbf{Exposición:} P \(\times\) I.
  \item \textbf{Severidad:} crítico (6–9), significativo (3–4), moderado (1–2); en la tabla se abrevian como Crít., Sign. y Mod.
  \item \textbf{Estado:} situación del riesgo tras la última revisión; los estados posibles (cap.~7 de la memoria) son identificado, sin novedad, activo, materializado, resuelto y cerrado.
\end{itemize}

% Contenido literal de la tabla 7.2 de la memoria (v0), más la columna
% Estado (novedad de las versiones: la v0 no la necesita porque sus estados
% viven en la tabla de seguimiento del capítulo). \footnotesize como en la
% v0. Anchos para 16 cm: 12,15 cm de columnas p{} + ~0,8 cm de las tres
% columnas c + separadores/filetes con \tabcolsep en 4 pt.
{\footnotesize
\setlength{\tabcolsep}{4pt}
% \hspace{0pt} al abrir cada celda: TeX no guionea nunca la PRIMERA palabra
% de un párrafo (y cada celda lo es); sin ese hueco de anchura cero,
% «Identificado» no puede partirse y desborda la columna Estado.
\begin{longtable}{|>{\RaggedRight\arraybackslash\hspace{0pt}}p{0.5cm}|>{\RaggedRight\arraybackslash\hspace{0pt}}p{3.9cm}|c|c|c|>{\RaggedRight\arraybackslash\hspace{0pt}}p{0.75cm}|>{\RaggedRight\arraybackslash\hspace{0pt}}p{3cm}|>{\RaggedRight\arraybackslash\hspace{0pt}}p{2.7cm}|>{\RaggedRight\arraybackslash\hspace{0pt}}p{1.3cm}|}
\caption{Registro de riesgos (versión 1)} \label{tab:plan-v1-registro} \\
\hline
\textbf{Id} & \textbf{Riesgo} & \textbf{P} & \textbf{I} & \textbf{Exp.} & \textbf{Sev.} & \textbf{Indicador clave} & \textbf{Indicador de resolución} & \textbf{Estado} \\ \hline
\endfirsthead

\multicolumn{9}{c}{{\bfseries \tablename\ \thetable{} — continuación de la página anterior}} \\
\hline
\textbf{Id} & \textbf{Riesgo} & \textbf{P} & \textbf{I} & \textbf{Exp.} & \textbf{Sev.} & \textbf{Indicador clave} & \textbf{Indicador de resolución} & \textbf{Estado} \\ \hline
\endhead

\hline \multicolumn{9}{|r|}{{Continúa en la página siguiente...}} \\ \hline
\endfoot

\hline
\endlastfoot

R1 & \textbf{Crecimiento descontrolado del alcance.} Producto ambicioso (\eng{feed}, mapa, chat, IA) con el autor como cliente y desarrollador a la vez: nadie externo frena la incorporación de funciones no previstas. & 3 & 3 & 9 & Crít. & El \eng{sprint} activo gana más de 2 puntos de historia tras la planificación del lunes sin retirar un volumen equivalente, o se implementa una función ausente del capítulo 3 sin haber pasado por el \eng{backlog}. & — & Identificado. \\ \hline

R2 & \textbf{Error sistemático de estimación del esfuerzo.} Primera estimación en solitario con puntos de historia y sin velocidad histórica propia: la hipótesis inicial de equivalencia entre puntos y horas (capítulo 10) puede estar sesgada de raíz. & 3 & 2 & 6 & Crít. & La velocidad real de los tres primeros \eng{sprints} se desvía más de un 25\,\% de la hipótesis inicial. & — & Identificado. \\ \hline

R3 & \textbf{Memoria relegada frente al desarrollo.} La tentación natural de un TFG con producto es programar ahora y documentar después, perdiendo el porqué de las decisiones que alimenta esta memoria. & 2 & 3 & 6 & Crít. & Un capítulo afectado por trabajo ya cerrado acumula dos \eng{sprints} sin actualizarse, o el diario de una sesión se redacta más de 48 horas después de esta. & — & Identificado. \\ \hline

R4 & \textbf{Incumplimiento de la normativa de la memoria.} La guía de estilo es normativa: una afirmación externa sin cita se trata como plagio (suspenso directo). El volumen de fuentes y el uso de IA generativa hacen posible que se cuele texto sin referenciar. & 2 & 3 & 6 & Crít. & Una revisión de capítulo encuentra una afirmación no original sin su cita. & La revisión final previa al depósito concluye con cero afirmaciones sin cita, con auditoría manual junto con ayuda de IA. & Identificado. \\ \hline

R5 & \textbf{Dependencia acrítica de la IA generativa.} El desarrollo asistido por IA es parte del método de trabajo; el riesgo es integrar código o texto generado sin la comprensión necesaria para defenderlo ante el tribunal, o con defectos de calidad o seguridad que una lectura superficial no detecta. & 2 & 3 & 6 & Crít. & Al redactar el diario, una decisión o un bloque integrados esa sesión no pueden explicarse sin recurrir de nuevo a la herramienta, o una revisión encuentra código integrado sin revisar ni probar. & — & Identificado. \\ \hline

R6 & \textbf{Curva de aprendizaje de la pila tecnológica.} PostGIS, pgvector, Prisma, la mensajería y las actualizaciones en tiempo real, las colas de trabajos en segundo plano y las notificaciones \eng{push} son territorio nuevo para el autor: buena parte del servidor se aprende sobre la marcha, con potencial de consumir días en un problema imprevisto. & 3 & 2 & 6 & Crít. & Una historia o una investigación previa consumen el doble de su tiempo previsto sin resolución a la vista. & — & Identificado. \\ \hline

R7 & \textbf{Incumplimiento del RGPD en el diseño.} El producto maneja ubicación de personas, verificación de identidad con documento y la regla de mayoría de edad: datos especialmente delicados. El reglamento es además estricto y transversal: casi cada función del producto toca datos personales, sobre todo al diseñar el modelo de datos, de modo que la ocasión de error se repite durante todo el desarrollo. & 2 & 3 & 6 & Crít. & Una entidad del modelo de datos incorpora un dato personal sin finalidad declarada, o un requisito exige almacenar la imagen del documento de identidad. & El modelo de datos completo pasa la revisión de minimización: cada dato personal con finalidad declarada y ningún documento de identidad almacenado. & Identificado. \\ \hline

R8 & \textbf{Capacidad real inferior a la declarada.} La planificación (capítulo 10) ya incorpora las restricciones conocidas —los compromisos laborales del verano, la estancia Erasmus del primer cuatrimestre y las prácticas de empresa desde febrero—, pero las horas declaradas por periodo siguen siendo hipótesis: la capacidad real puede quedar por debajo incluso de esa previsión reducida. & 2 & 2 & 4 & Sign. & Las horas reales registradas en el diario quedan por debajo del 70\,\% de las declaradas durante dos semanas consecutivas sin causa puntual identificada. & — & Identificado. \\ \hline

R9 & \textbf{Dependencia de servicios externos.} La IA de plazos, la base de datos gestionada y la verificación de identidad dependen de terceros con costes, cuotas y disponibilidad ajenos al proyecto; además, el proyecto se autoimpone no pagar por ningún servicio, de modo que solo valen las capas gratuitas. & 2 & 2 & 4 & Sign. & Anuncio de cambio de precios o de límites de un servicio, agotamiento de su capa gratuita, o indisponibilidad superior a 24 horas. & El núcleo del producto funciona de punta a punta con versiones simuladas de los servicios externos, sin que ninguna prueba requiera uno real al finalizar la implementación de la aplicación y la versión final de la memoria. & Identificado. \\ \hline

R10 & \textbf{Indisponibilidad prolongada del autor.} Una enfermedad o imprevisto personal detiene por completo un proyecto unipersonal: no hay a quién traspasar el trabajo. & 1 & 3 & 3 & Sign. & Más de tres días consecutivos sin poder trabajar por causa sobrevenida dentro de un periodo de dedicación plena. & — & \textbf{Resuelto} (materializado en S1, superado el 26-07-2026). \\ \hline

R11 & \textbf{Pérdida de trabajo o del entorno local.} Un fallo del único equipo de desarrollo o un borrado accidental puede destruir trabajo no respaldado. & 1 & 3 & 3 & Sign. & Reinicios espontáneos, errores de disco o arranques anómalos del equipo, o más de dos sesiones consecutivas sin subir los cambios al remoto. & — & Identificado. \\ \hline

R12 & \textbf{Ciclos de revisión con el tutor.} La memoria depende de revisiones externas: una respuesta que se demora o una revisión que obliga a rehacer un capítulo dado por cerrado retrasan el cierre. La relación previa con el tutor y el cauce de comunicación ya establecido mantienen la probabilidad baja. & 1 & 2 & 2 & Mod. & Una consulta o un capítulo enviado a revisión acumula más de dos semanas sin respuesta, o una revisión devuelve cambios estructurales sobre un capítulo dado por cerrado. & El borrador completo de la memoria pasa una revisión del tutor sin cambios estructurales. & Identificado. \\ \hline

R13 & \textbf{Fricción imprevista de la cadena de herramientas.} Las fricciones ya conocidas del entorno Windows (compilación LaTeX sin Perl, enlaces simbólicos de pnpm) están resueltas y documentadas: son problemas gestionados, no riesgos. El riesgo es que la combinación Windows/Docker/Expo/MiKTeX produzca incompatibilidades nuevas que consuman capacidad sin aportar valor. & 2 & 1 & 2 & Mod. & Más de una hora de una misma sesión perdida en un problema de entorno ajeno al producto y sin receta previa. & — & Identificado. \\ \hline

\end{longtable}
}

\section{Plan de respuesta ante riesgos}

% Contenido literal de la tabla 7.3 de la memoria (v0), con el único cambio
% documentado en la NOTA (R4: mitigación sin la cláusula de páginas).
% \small como en la v0. Anchos del .tex de referencia: 0,8 + 6,9 + 6,9 cm
% + separadores = 15,9 cm, dentro de los 16 cm de texto.
{\small
\begin{longtable}{|>{\RaggedRight\arraybackslash}p{0.8cm}|>{\RaggedRight\arraybackslash}p{6.9cm}|>{\RaggedRight\arraybackslash}p{6.9cm}|}
\caption{Plan de respuesta (versión 1)} \label{tab:plan-v1-respuesta} \\
\hline
\textbf{Id} & \textbf{Mitigación (prevención)} & \textbf{Contingencia} \\ \hline
\endfirsthead

\multicolumn{3}{c}{{\bfseries \tablename\ \thetable{} — continuación de la página anterior}} \\
\hline
\textbf{Id} & \textbf{Mitigación (prevención)} & \textbf{Contingencia} \\ \hline
\endhead

\hline \multicolumn{3}{|r|}{{Continúa en la página siguiente...}} \\ \hline
\endfoot

\hline
\endlastfoot

R1 & El alcance queda delimitado por contrato en el capítulo 3 (qué entra y qué queda fuera), cada historia del \eng{backlog} está marcada como imprescindible, deseable u opcional (capítulo 10), y toda idea nueva entra por el \eng{backlog}, nunca directamente al \eng{sprint} en curso. & Descartar las funciones opcionales, renegociar las deseables y mover lo recortado a líneas futuras. \\ \hline

R2 & Medir la velocidad real desde el primer \eng{sprint}, recalibrar la equivalencia horas/punto en cada cierre y trocear cualquier historia grande antes de abordarla. & Recortar lo opcional para proteger lo imprescindible y ampliar la horquilla de cierre, dejando constancia de la desviación en el capítulo 13. \\ \hline

R3 & La redacción tiene hueco reservado en todos los \eng{sprints} (capítulo 10), ninguna historia se da por cerrada hasta que su decisión queda documentada, y el diario de desarrollo se escribe al cierre de cada sesión, no de memoria semanas después. & Dedicar \eng{sprints} exclusivos a la memoria antes del cierre (la última fase ya lo prevé parcialmente). \\ \hline

R4 & Toda afirmación procedente de una fuente externa se cita en el momento de escribirla, nunca en una pasada posterior; las entradas de bibliografía siguen las plantillas del repositorio; los pendientes se marcan con nota visible y se purgan antes del depósito. & Auditoría completa de citas, capítulo a capítulo contra sus fuentes, antes del depósito. \\ \hline

R5 & Regla ya operativa del proyecto: cada decisión se explica y justifica en el diario con su alternativa descartada, todo código generado se revisa antes de integrarse y nada se da por terminado sin pasar las comprobaciones automáticas de calidad y las pruebas. & Congelar el módulo afectado, estudiarlo y reescribirlo hasta poder explicarlo sin asistencia, y registrar el episodio en el diario como material para la reflexión sobre el uso de IA. \\ \hline

R6 & Antes de comprometer historias que dependan de una tecnología desconocida, dedicarle una investigación previa con límite de tiempo explícito, y preferir servicios gestionados donde exista esa opción. & Adoptar la alternativa documentada en el análisis de alternativas o recortar la función a su versión mínima. \\ \hline

R7 & Minimización desde el diseño, ya decidida en las reglas de producto, el alcance y los requisitos (capítulos 2, 3 y 5): la verificación se delega en un proveedor externo guardando solo el estado y una clave de referencia (\eng{token}), la ubicación se difumina y exige consentimiento, y los datos residen en región europea. & Recortar o desactivar la función afectada (por ejemplo, el mapa de descubrimiento) hasta rediseñarla conforme al reglamento. \\ \hline

R8 & Horas reales por sesión en el diario; cada retrospectiva las compara con las declaradas; calendario con holgura hasta la defensa (septiembre como reserva). & Rebajar la capacidad declarada de los \eng{sprints} restantes al nivel observado; recortar primero lo opcional. \\ \hline

R9 & Todas las integraciones externas se declaran objetivos secundarios o se desarrollan contra versiones simuladas que funcionan sin el servicio real, de forma que el núcleo del producto nunca dependa de ellas; se eligen proveedores con capa gratuita suficiente para un uso académico. & Sustituir el servicio por una versión simulada o por una alternativa gratuita equivalente, y documentar la integración real como línea futura. \\ \hline

R10 & Margen temporal deliberado entre el cierre objetivo y la defensa (la convocatoria de septiembre de \thesisYear{} actúa como reserva) y documentación continua (diario más memoria) que permite retomar el contexto tras una pausa sin pérdida. & Pausa formal, re-planificación del calendario y, en el extremo, cambio de convocatoria. \\ \hline

R11 & Absolutamente todo (código, memoria, diario, planificación) vive en el mismo repositorio git con remoto en GitHub, con subidas frecuentes; esta es precisamente una de las razones por las que la documentación del proyecto se versiona junto al código. & Restaurar desde el remoto; la pérdida máxima queda acotada a la sesión en curso. \\ \hline

R12 & Enviar los capítulos a revisión conforme se van cerrando, en lugar del borrador completo al final, y acordar en cada envío un plazo orientativo de respuesta. & Continuar con trabajo que no dependa de la revisión pendiente y re-planificar el capítulo afectado cuando llegue la respuesta. \\ \hline

R13 & Cada solución encontrada se documenta como receta reproducible en el repositorio en el momento de resolverla, para pagar cada problema una sola vez. & Aplicar la alternativa documentada (por ejemplo, el modo alternativo de instalación de dependencias de pnpm) o trasladar el paso conflictivo a la integración continua, que corre en Linux. \\ \hline

\end{longtable}
}

\end{document}
